% SPDX-License-Identifier: Apache-2.0
\section{Whether a Receive Took Something}
\label{sec:x-recvsome}

\code{recv\_if(c)} takes a transaction when \code{c} holds and the
channel offers one, and gives it back as \code{Some}. Asked
\code{is\_some()}, the answer is that take: one bit, the condition and
the channel's valid. A unit names the receive, \code{let got =
inp.recv\_if(go)}, and asks it afterwards; the transaction itself is
\code{got}'s data, and \code{got.is\_some()} is not.

\code{Takers} holds a \code{Source}, which offers its next count
whenever the channel has room, and a \code{Taker}, which takes a word
on every other cycle, shows on \code{took} whether it took one, and
counts what it took. The channel transfers a word and a flag, nine bits,
so a lowering that gave the take as the data would be wrong in a way
both netlists show. The lowering once did (issue 1079): nvc refused the
VHDL, a nine-bit value where a bit belongs, and Verilator ran the
Verilog, where \code{took} showed the data's low bit and the count went
up for any word on the channel that was not zero, taken or not. The
run asserts eight takes, and both netlists are checked against it.

\begin{figure}[htbp]
\centering
\input{recv_some_timing.tex}
\caption{A take on every other cycle: \code{took} is high each time
\code{go} is and a word is offered, and \code{count} climbs by one at
each, whatever the word.}
\label{fig:x-recvsome}
\end{figure}

\lstinputlisting[caption={\code{ex\_recv\_some.rs}. Run with \code{bazel run //lib/examples:ex\_recv\_some}.},label={lst:x-recvsome}]{lib/examples/ex_recv_some.rs}

\lstinputlisting[style=ascii,caption={What \code{ex\_recv\_some} printed when the build ran it: the takes counted, and the netlist, with the taker's \code{took} driven by the channel's valid and the condition.},label={lst:x-recvsome-out}]{out_recv_some.txt}
